% explainer-source-hash: 960ce50d1a06f5cb60d0eb961093288892eeda38be15723bf3830b5550d3321c
% explainer-source-commit: 3f6c7a8913b4ee27f99ff1a17bb5442bba11de07
% explainer-generated: 2026-07-16
% Regenerate with /explain-all (see WIRING.md). Do not edit this stamp by hand:
% it is how the toolchain knows whether this document still matches the code.
\documentclass[11pt,a4paper]{article}
\input{preamble}

\title{\vspace{-1.5cm}\textbf{Dino Game Bot}\\[0.3cm]
\large A Brief: What It Does, How It Decides, and Where It Is Soft}
\author{Portfolio explainer: concise overview}
\date{}

\begin{document}
\maketitle
\thispagestyle{fancy}

\begin{keyidea}{The project in four sentences}
Chrome hides an endless-runner game at \code{chrome://dino/}: a dinosaur runs, cacti
approach, you press space to jump. This project plays it without touching the game's
code or data, by screenshotting one rectangle of the screen and asking whether any
pixel in it is the colour an obstacle would be. If one is, it sends a space keypress.
812 lines of Python, three files, three dependencies.
\end{keyidea}

\tableofcontents
\newpage

\section{The shape of the thing}

Everything the bot does hangs off one loop, and the loop is four steps long.

\begin{figure}[H]
\centering
\resizebox{0.98\textwidth}{!}{
\begin{tikzpicture}[node distance=6mm and 13mm]
  \node[box] (grab1) {\textbf{1. Which mode?}\\screenshot a large box,\\count white vs. grey pixels};
  \node[box, right=of grab1] (grab2) {\textbf{2. Anything there?}\\screenshot a small box,\\look for the obstacle colour};
  \node[dec, right=of grab2] (d) {found?};
  \node[box, right=of d] (jump) {\textbf{3. Jump}\\send SPACE\\to the browser};
  \node[box, below=11mm of d] (rec) {\textbf{4. Tell the dashboard}\\mode, obstacle, jump count};

  \draw[flow] (grab1) -- node[lbl, above] {mode} (grab2);
  \draw[flow] (grab2) -- (d);
  \draw[flow] (d) -- node[lbl, above] {yes} (jump);
  \draw[flow] (jump) |- (rec);
  \draw[flow] (d) -- node[lbl, right] {no} (rec);
  \draw[dflow] (rec.west) -- ++(-42mm,0) |- (grab1.south);
  \node[lbl, below=1mm of grab1.south east, xshift=6mm] {repeat};
\end{tikzpicture}
}
\caption{The whole bot. Step 1 exists only because step 2's answer depends on it.}
\end{figure}

\section{Why step 1 exists}

The Dino game inverts its palette periodically: day mode is dark sprites on a white
background, night mode is white sprites on a dark background. So ``what colour is an
obstacle?'' has two answers, and the wrong one is not merely inaccurate, it is
inverted: a bot hunting for dark pixels in night mode would find the background
everywhere and jump forever.

Hence the sequence. \code{check\_day} takes the majority vote of white against dark
grey across a large region, which settles the mode; \code{check\_obstacle} then hunts
for whichever colour the \emph{other} mode uses, inside a much smaller region placed
just in front of the dinosaur.

\begin{keyidea}{The intelligence is in the rectangle, not the code}
\code{check\_obstacle} does not measure distance, speed, or shape. It asks whether a
single pixel of the target colour appears anywhere in its box. ``Obstacle in the
box'' and ``jump now'' are treated as the same event, which is only true because the
box was hand-placed at exactly the right spot. That makes the four hardcoded
coordinates the most load-bearing and most fragile thing in the repository.
\end{keyidea}

\section{The three good decisions}

Worth knowing, because each is defensible in an interview.

\subsection{Demo mode patches the eye, not the brain}

The bot needs a live Chrome playing a live game on a screen matching its hardcoded
coordinates. A fresh clone has none of that, so \code{--demo} replays six pre-drawn
frames instead. The interesting part is \emph{how}.

\begin{lstlisting}[language=Python,caption={The substitution, from \code{run\_demo}}]
original_grab = ImageGrab.grab
try:
    for filename, frame in frames:
        ImageGrab.grab = lambda bbox=None, _frame=frame: _frame.crop(bbox)
        is_day = check_day()
        obstacle = check_obstacle(is_day)
        ...
finally:
    ImageGrab.grab = original_grab
\end{lstlisting}

Rather than adding an \code{if demo:} branch inside the detection functions, it
swaps out the screen-capture call those functions make, \code{ImageGrab.grab}, for
a stand-in that crops a stored image. The detection functions are byte-for-byte
the ones live mode runs, and contain no code that could tell the difference.

\begin{keyidea}{Why that distinction is the whole point}
A demo-only branch means demo mode tests a path production never takes, which is how
tests come to pass while the product is broken. Patching the input leaves the logic
untouched, so what runs really is what ships. The \code{finally} restore matters too:
without it the process would keep returning a stale dino frame from any later capture.
\end{keyidea}

\subsection{The fixture imports the constants it tests against}

\code{tools/generate\_demo\_frames.py} could have hardcoded the colours and
coordinates when drawing its frames. It imports them from \file{main.py} instead, so
re-measuring the coordinates for a different screen moves the test frames with them.
A fixture that duplicates the constants of the code it tests is a fixture that will
eventually disagree with it silently.

\subsection{A change was only kept because it could be proven}

\code{check\_day} and \code{check\_obstacle} gained an optional \code{tolerance}
argument: instead of requiring a pixel to be bit-for-bit \code{(255,255,255)}, count
it if every channel is within \code{tolerance}. This matters because real screen
captures drift (anti-aliasing, dithering, compression) in ways a hand-drawn flat
colour never does.

\code{tools/compare\_detection\_robustness.py} tests it properly: both logics against
both a clean and a deliberately noisy frame set, so ``did it help?'' and ``did it
break what worked?'' are answered by the same run, with the expected answers derived
from filenames rather than a hand-maintained table.

\begin{table}[H]
\centering
\begin{tabular}{lcc}
\toprule
 & \textbf{clean frames} & \textbf{noisy frames} \\
\midrule
old (\code{tolerance=0}) & 6/6 correct & 4/6 correct \\
new (\code{tolerance=8}) & 6/6 correct & 6/6 correct \\
\bottomrule
\end{tabular}
\caption{I ran this. The README's table reproduces exactly, including which two
frames the old logic fails.}
\end{table}

And it defaults to \code{0}, the original exact-match behaviour, because it has only
ever been proven against synthetic noise. Offering a change rather than imposing it,
on the grounds that the evidence is synthetic, is the right instinct.

\section{Where it is soft}

The findings below come from running the code, not reading it.

\begin{honest}{The noise in the benchmark only applies to one of the two colours}
\code{dither()} shifts every channel by $\pm 6$ and its docstring promises ``no pixel
in the output is ever exactly equal to a color present in the input''. It is not
true. Dark grey 83 moves to 89 or 77, both genuinely off-target. White 255 moves to
249 or to \code{min(255, 261)}, which is 255 again. Measured: 420,000 pixels of the
noisy \file{01\_day\_clear.png}, exactly half, come through bit-for-bit white. So the
noise fully perturbs grey and half-ignores white, and the experiment is not measuring
what it thinks. Fixing it is one line, and would likely make the tolerance feature
look better rather than worse, but as written the evidence does not support the
conclusion it is offered for.
\end{honest}

\begin{honest}{A tie in \code{check\_day} silently votes ``night''}
The return is \code{day\_count > night\_count}. A frame containing neither colour
scores $0 > 0$, which is \code{False}, so the function reports night on zero
evidence. Verified against a uniform grey frame. This is why two of the old logic's
four ``correct'' answers on the noisy set are worthless: \file{04} and \file{05} are
night frames, the old logic saw nothing at all, defaulted to night, and was scored
correct. Its honest score on that set is closer to 2/6 than 4/6. The benchmark cannot
distinguish ``right'' from ``no opinion, guessed, got lucky''.
\end{honest}

\begin{honest}{The recommended default would slow the bot roughly ninefold}
The README's ``What I'd do differently'' proposes turning tolerance on by default.
Nothing anywhere measures what that costs. I did, on this machine, excluding the real
screen grab: one loop iteration's pixel work goes from about \textbf{70 ms} at
\code{tolerance=0} to about \textbf{608 ms} at \code{tolerance=8}, roughly 14 Hz down
to 1.6 Hz. The cause is that \code{count()} and \code{in} are C-speed list scans,
while the tolerant path is a Python generator calling \code{\_color\_distance} once
per pixel, about 140,000 function calls where there was one C loop. A bot forming an
opinion 1.6 times a second will hit obstacles while still counting pixels from a
stale frame. The feature is correct; accuracy was simply the only axis measured, and
for a real-time loop it is not the deciding one. It is fixable (do the comparison in
C via Pillow's \code{point()}, or sample a sparse grid instead of 140,400 pixels),
but ``turn it on by default'' as stated would break the bot.
\end{honest}

\begin{honest}{The advertised 10 ms polling loop actually runs near 14 Hz}
The README frames the 10 ms sleep as the reaction-time knob. It is not: it sits next
to roughly 70 ms of unavoidable pixel scanning, and it only executes on iterations
that jumped. On a clear track the loop never sleeps at all, pinning one core at 100\%
while it takes two screenshots and scans 206,000 pixels per pass. The real lever is
the pixel count. \code{check\_day} scans 140,400 pixels on every iteration to answer
one binary question that a few hundred would answer, about a palette that flips maybe
once a minute. Caching the mode would remove roughly 70\% of the loop's cost for no
loss.
\end{honest}

\begin{honest}{The demo frames omit the feature most likely to break the bot}
\code{make\_frame} fills the canvas with one flat colour, and says so: ``the real
game has sky and a thin ground strip''. That simplification is harmless for
\code{check\_day}, which counts and compares. It is not harmless for
\code{check\_obstacle}, which fires if the target colour appears \emph{once} in its
box. In the real day-mode game the ground line is dark grey, the exact colour
\code{check\_obstacle} hunts for in day mode, and the box spans the vertical band
where a ground line would plausibly sit. If it falls inside, the bot jumps
continuously forever. To be precise: that the frames contain no ground strip is
\textbf{confirmed} (their histograms hold two colours, no third); that the real
game's ground line intersects the box is \textbf{inferred} and untestable here, with
no Chrome and no live game. The point survives either way: the fixture cannot tell
those two worlds apart, so ``six of six pass'' proves less than it sounds. The
fixture was drawn by the same person holding the same assumption as the code.
\end{honest}

\begin{honest}{Smaller items}
The constants named \code{DAY\_GROUND\_COLOR}/\code{NIGHT\_GROUND\_COLOR} hold
background colours, not ground colours, and the docstrings inherit the confusion.
On a headless Linux box that is not running as root, there is no working way to stop
the loop: the console view has no \code{closed} attribute, and the \code{q} hotkey
needs root to read the input device. \code{driver.find\_element(value="t")} is
re-queried on every jump instead of hoisted. And two of the README's eight
``Challenges'' bullets are about LaTeX errors in the explainer PDFs, which are not
challenges of this project at all.
\end{honest}

\section{Verified against not verified}

\begin{figure}[H]
\centering
\begin{tikzpicture}[node distance=4mm and 10mm]
  \node[box, fill=okgreen!10, draw=okgreen, text width=52mm] (v)
    {\textbf{Verified here}\\[2pt]
     \scriptsize
     demo runs, 6/6 correct decisions\\
     headless fallback prints the same\\
     \texttt{ImageGrab.grab} restored after\\
     benchmark table reproduces exactly\\
     frame histograms match the geometry\\
     tie in \texttt{check\_day} votes night\\
     \texttt{dither()} leaves white intact\\
     the 70 ms / 608 ms timings};
  \node[extbox, right=of v, draw=warnred, text width=52mm] (n)
    {\textbf{Not verified, needs real Chrome}\\[2pt]
     \scriptsize
     that the game loads and responds\\
     that the coordinates match any\\ real window today\\
     whether the ground line breaks\\ obstacle detection\\
     whether it clears obstacles at all\\
     whether tolerance helps against\\ real capture noise};
  \draw[flow] (v) -- (n);
\end{tikzpicture}
\caption{The line this project cannot cross without a machine running the game. The
README draws the same line honestly; this document adds the measurements.}
\end{figure}

\section{Glossary}

\begin{description}[leftmargin=0pt, style=nextline]
\item[RGB] A colour as three numbers 0 to 255: red, green, blue. White is
  \code{(255,255,255)}; the game's dark grey is \code{(83,83,83)}.
\item[Bounding box] A screen rectangle as \code{(left, top, right, bottom)} in
  pixels from the top-left corner. Down is positive.
\item[Screen scraping] Driving a program by reading its pixels and sending fake
  input, because it offers no API. Fragile by nature: the only contract is that the
  picture keeps looking the same.
\item[Selenium / WebDriver] A library that drives a real browser, through a helper
  program (\code{chromedriver}). Since 4.6 it fetches that helper itself, which is
  why no manual driver setup is needed.
\item[Pillow / \code{ImageGrab}] Python's imaging library; \code{ImageGrab.grab(bbox)}
  captures a screen rectangle. It is the bot's eye and the one thing demo mode
  replaces.
\item[Monkeypatching] Replacing a module's function at runtime from outside it.
  Demo mode's mechanism.
\item[tkinter / \code{mainloop()}] Python's built-in windowing toolkit.
  \code{mainloop()} never returns, so instead the dashboard pumps the event queue
  once per iteration and returns immediately, keeping the window alive without
  taking the loop's thread.
\item[\code{deque(maxlen=n)}] A queue that discards its oldest item automatically
  when full. Makes ``keep the last 8 events'' a data-structure choice.
\item[Chebyshev distance] Difference between two colours measured as the largest
  gap on any single channel. What \code{tolerance} compares against.
\end{description}

\begin{keyidea}{The summary}
A small project that does one thing, unusually well-instrumented for its size, with
three genuinely good instincts: exercise production code from the test path, keep one
source of truth for the constants, and refuse a change the evidence does not support.
Its soft spot is that the evidence it does trust measures accuracy on frames whose
noise is half-broken, never measures speed at all, and cannot see the failure mode
most likely to be real. All of that is cheap to fix, and much better learned here
than from an interviewer.
\end{keyidea}

\end{document}
